\section{Conclusion}

\subsection{Summary}

This study developed a hybrid monitoring scheme in which a CAE supplies a reconstruction error and
learned features, an OC-SVM scores those features, and a MEWMA chart accumulates the vector
$(s_t,\ell_t)$. It was implemented on a simulated nonlinear process with five independent datasets,
and every chart was calibrated empirically to an in-control ARL of 370. The learned components
behaved largely as intended, and on independent in-control data the hybrid was somewhat more liberal,
and the classical benchmark somewhat more conservative, than the target. The hybrid had a lower
$\ARL_1$ than both single-component ablations at every shift (within sampling error only against
$S(0.05)$ at $\delta\le0.4$), so the combined vector carries more usable information than either
learned signal alone, but it did not outperform the classical MEWMA: the two were
indistinguishable at $\delta=0.2$, and the benchmark was faster at every larger shift, by up to 106
monitoring opportunities on average, with the hybrid needing 1.4 to 2.5 times as many post-change
observations. Larger smoothing constants shortened the hybrid's run lengths modestly, most in
relative terms at large shifts, without closing the gap.

\subsection{Recommendations}

First, for sustained mean shifts in a moderately nonlinear process, a well-calibrated classical MEWMA
on the raw variables should remain the default, and a learned monitor should be adopted only where it
improves run-length performance against such a benchmark under a common in-control target. Second,
learned reconstruction and one-class information should be monitored jointly rather than separately,
since the joint chart dominated both ablations. Third, charts built on overlapping windows should be
calibrated empirically on the dependent process, with delays reported both in monitoring
opportunities and in post-change observations.

\subsection{Future work}

The hybrid should be evaluated on faults that alter the shape, variability or dependence of the data
(variance changes, sensor faults, drifts and transients), for which distance-type features may be
better suited than for a mean shift, and on larger, more strongly nonlinear processes.
Direction-preserving learned features, dependence-aware finite-start scaling for overlapping windows,
changes that occur during monitoring, and validation on the Tennessee Eastman Process and real plant
data also deserve study.
